% Transcription — fonds Alexandre Grothendieck, shelfmark 19, batch 1 (pages 1-20).
% Established from the facsimile archives/batches/19/batch-01.pdf (a mirror of
% https://grothendieck.umontpellier.fr/19.pdf), per the skill
% transcribe-grothendieck.
%
% Pages 2 and 12 are blank versos — absent.
%
% The batch holds two distinct pieces, and the folder cover (page 1) names both:
% "Résolutions standards et foncteurs adjoints" and "Tapis cartésiens", with an
% archivist's pencilled tally "9p + 14p" beside them. The tally matches what is
% here. Pages 3-11 are the nine sheets of "Tapis cartésiens", numbered 1) 2) 3)
% by the author: the comonad φ = fg attached to an adjoint pair, Godement's
% conditions, the category K(φ,α,λ) of coalgebras, a monadicity theorem, and the
% case B = ∏ B_i worked down to two factors. Pages 13-20 open the second piece —
% one sheet headed "Résolutions standard et théories relatives", then a text
% headed "Cohomologie relative" and numbered 1., 2., 3. — Yoneda, S-exactness,
% S-injectives, a canonical S-resolution. Fourteen pages of it would run to page
% 26, i.e. past this batch and into batch 2.
%
% The dating is the inventory's own, for the whole shelfmark; nothing on these
% pages carries a date.
%
% Pass: Opus 5 (claude-opus-5), 2026-08-22 — first pass, unchecked against the pages by a human.
\documentclass[11pt,a4paper]{article}
% Le préambule commun à toutes les transcriptions du fonds Grothendieck.
%
% Il tient en une idée : **l'incertitude se compose, elle ne se résout pas.**
% Une transcription qui lisse un mot douteux a détruit la seule chose pour
% laquelle on la faisait — un lecteur doit pouvoir savoir, à chaque endroit,
% ce qui était sur le feuillet et ce qui a été ajouté. Les six macros qui
% suivent sont donc l'appareil critique, et non de la décoration.
%
% Les mêmes commandes sont reconnues par `scripts/render.mjs`, qui produit la
% vue de lecture du site. Une macro ajoutée ici doit l'être aussi là-bas,
% faute de quoi le rendu échouera bruyamment — ce qui est voulu : un rendu
% silencieusement faux est pire qu'un rendu absent.

\ProvidesPackage{grothendieck}[2026/08/08 Transcriptions du fonds Grothendieck]

\usepackage{amsmath,amssymb,amsthm}
\usepackage{mathrsfs}
\usepackage{stmaryrd}
% Les diagrammes commutatifs sont partout dans ce fonds : c'est la langue
% même de la géométrie algébrique de Grothendieck, et une transcription qui
% les rendrait en prose ne transcrirait rien.
\usepackage{tikz-cd}
% Français par défaut — c'est la langue des feuillets — mais l'anglais est
% chargé aussi : la traduction et le résumé sont des documents anglais, et la
% typographie française (espaces avant les deux-points) y serait une faute.
% Ces documents basculent par \selectlanguage{english} en tête de corps.
\usepackage[english,french]{babel}
\usepackage{fontspec}
\usepackage{geometry}
\geometry{a4paper,margin=25mm,marginparwidth=28mm}
\usepackage{marginnote}
\usepackage{xcolor}
% ulem, not soul. `soul`'s \st and \ul are built on a character-by-character
% scan that XeLaTeX's font machinery breaks: the document fails with an
% undefined control sequence rather than degrading. ulem does the same two jobs
% and is Unicode-engine safe. `normalem` stops it hijacking \emph, which the
% transcriptions use for Grothendieck's own emphasis.
\usepackage[normalem]{ulem}
\usepackage{hyperref}
\hypersetup{colorlinks=true,linkcolor=black,urlcolor=black,pdfborder={0 0 0}}

% --- Métadonnées du lot -----------------------------------------------------
% Reprises telles quelles par le script de rendu, qui les affiche en tête de
% la vue de lecture. Elles fixent ce que le fichier prétend couvrir : sans
% elles, une transcription flotte sans point d'attache dans un fonds de seize
% mille feuillets.

\newcommand{\folder}[1]{\gdef\grothFolder{#1}}
\newcommand{\batch}[1]{\gdef\grothBatch{#1}}
\newcommand{\leaves}[2]{\gdef\grothFirst{#1}\gdef\grothLast{#2}}
\newcommand{\foldertitle}[1]{\gdef\grothFoldertitle{#1}}
\gdef\grothFolder{?}\gdef\grothBatch{?}\gdef\grothFirst{?}\gdef\grothLast{?}\gdef\grothFoldertitle{}

% --- Le résumé --------------------------------------------------------------
%
% Il ouvre la lecture modernisée, sous le titre « Résumé », et il est composé
% en petit. Ce n'est pas un ornement : c'est le seul passage du document qui
% ne soit pas des mathématiques, et un lecteur doit pouvoir le reconnaître
% avant de l'avoir lu — puis le sauter s'il connaît déjà le sujet, ou s'y
% arrêter s'il ne le connaît pas. Le filet qui le suit marque la reprise du
% corps.

\newenvironment{resume}
  {\small\setlength{\parskip}{0.45em}}
  {\par\medskip\hrule\medskip}

% --- Mots-clés ---------------------------------------------------------------
%
% En anglais, en clôture du résumé de la lecture modernisée : le vocabulaire
% moderne sous lequel chercher ce que les feuillets construisent. Le manifeste
% du site (`npm run manifest`) les extrait pour étiqueter la cote entière —
% cette ligne est la source unique des tags, il n'y a pas de fichier à côté.
\newcommand{\keywords}[1]{\par\smallskip\noindent
  {\footnotesize\textsc{Keywords} --- \textit{#1}}\par}

% --- L'appareil critique ----------------------------------------------------

% Le numéro de feuillet, en marge. C'est l'ancre qui permet de régler un
% désaccord sur une formule en regardant *un* feuillet plutôt qu'une chemise
% de six cents. Le site s'en sert aussi pour tourner les pages du fac-similé
% quand on fait défiler la transcription.
\newcommand{\leaf}[1]{%
  \marginnote{\footnotesize\color{gray!70}\textsf{#1}}%
  \label{leaf:#1}\ignorespaces}

% Illisible : on ne devine pas. Un mot inventé qui se lit comme les autres est
% le pire résultat possible.
\newcommand{\ill}{\textcolor{red!60!black}{[\,\ldots\,]}}

% Lecture douteuse : le mot est donné, et signalé comme incertain.
\newcommand{\uncertain}[1]{\uline{#1}}

% Ajout de l'éditeur — une lettre manquante, un mot sous-entendu.
\newcommand{\add}[1]{\textcolor{blue!50!black}{[#1]}}

% Biffé par Grothendieck lui-même. Ce qu'il a rayé fait partie du document :
% une rature dit souvent où la pensée a changé de direction.
\newcommand{\struck}[1]{\sout{#1}}

% Note du transcripteur, sur l'état matériel du feuillet.
\newcommand{\note}[1]{\par\smallskip{\small\color{gray!60!black}\textit{[#1]}}\par\smallskip}

% Note marginale de Grothendieck, distincte du corps du texte.
\newcommand{\marginal}[1]{\par\smallskip\noindent\hspace{1em}%
  {\small\color{gray!60!black}#1}\par\smallskip}

% --- Titre ------------------------------------------------------------------

\AtBeginDocument{%
  \begin{center}
    {\large\bfseries Fonds Alexandre Grothendieck --- cote n\textsuperscript{o}~\grothFolder}\\[2pt]
    {\itshape\grothFoldertitle}\\[4pt]
    {\small feuillets \grothFirst--\grothLast\ (lot \grothBatch)}\\[2pt]
    {\footnotesize\color{gray!60!black}Transcription établie d'après le fac-similé numérisé
     par l'Université de Montpellier.\\
     Les lectures incertaines sont soulignées, les illisibles notées [\,\ldots\,],
     les ajouts entre crochets.}
  \end{center}
  \vspace{1em}}

\endinput


\folder{19}
\batch{1}
\pages{1}{20}
\dating{[à partir de 1958-1973]}
\watermark{Édition de démonstration}
\foldertitle{Topos : notes manuscrites (s.d.)}

\begin{document}

\section*{Tapis cartésiens (pages 1 à 11)}

\page{1}
Titre porté à l'encre sur la chemise : « Résolutions standards et foncteurs
adjoints » ; au-dessous, « Tapis \uncertain{cartésiens} ».
\note{les deux titres correspondent aux deux textes du lot : « Tapis
cartésiens » occupe les pages 3 à 11, l'autre commence page 13}

\subsection*{1) Le foncteur $\varphi = fg$ et la catégorie $K(\varphi,\alpha,\lambda)$}

\page{3}
\emph{Tapis \uncertain{cartésiens}}

1) \quad
\begin{tikzcd}
A \arrow[r, bend left=20, "f"] & B \arrow[l, bend left=20, "g"]
\end{tikzcd}
\quad $g$ adjoint : \uncertain{droite} de $f$, d'où morphismes d'adjonction
\note{le mot « droite » est écrit par-dessus un premier mot raturé}
\[
  \mathrm{id}_{A} \xrightarrow{\ \beta\ } gf , \qquad
  fg \xrightarrow{\ \alpha\ } \mathrm{id}_{B}
  \tag{1.1}
\]
On pose
\[
  \varphi = fg : B \longrightarrow B
  \tag{1.2}
\]
On a donc
\[
  \varphi \xrightarrow{\ \alpha\ } \mathrm{id}_{B} , \qquad
  \varphi \xrightarrow{\ \lambda\ } \varphi^{2}
  \tag{1.3}
\]
i.e.
\[
  \lambda : fg \longrightarrow fgfg , \qquad \lambda = f * \beta * g
  \tag{1.4}
\]
Alors $\varphi$ satisfait aux \emph{conditions de Godement} :

(1.5) \quad (unitarité)

\begin{tikzcd}
\varphi \arrow[r, "\lambda"] \arrow[rr, bend right=35, "\mathrm{id}_{\varphi}"'] & \varphi^{2} \arrow[r, bend left=18, "\varphi*\alpha"] \arrow[r, bend right=18, "\alpha*\varphi"'] & \varphi
\end{tikzcd}

commutatif : $(\varphi * \alpha) \circ \lambda = \mathrm{id}_{\varphi}$.

(1.6) \quad (associativité)

\begin{tikzcd}
\varphi \arrow[r, "\lambda"] & \varphi^{2} \arrow[r, bend left=18, "\varphi*\lambda"] \arrow[r, bend right=18, "\lambda*\varphi"'] & \varphi^{3}
\end{tikzcd}

commutatif : $(\varphi * \lambda) \circ \lambda = (\lambda * \varphi) \circ \lambda$.

On définit la catégorie
\[
  K = K(\varphi, \alpha, \lambda)
  \tag{1.7}
\]
comme formée des couples
\[
  (X, u) \qquad X \in \mathrm{Ob}\, B, \quad u : X \longrightarrow \varphi(X)
  \tag{1.8}
\]
satisfaisant les conditions de commutativité

(1.10) \quad
\begin{tikzcd}
X \arrow[r, "u"] & \varphi(X) \arrow[r, bend left=18, "\varphi(u)"] \arrow[r, bend right=18, "\lambda(X)"'] & \varphi^{2}(X)
\end{tikzcd}
\quad $\lambda \circ u = \varphi(u) \circ u$ \quad (associativité)

(1.9) \quad
\begin{tikzcd}
X \arrow[r, "u"] & \varphi(X) \arrow[r, "\alpha(X)"] & X
\end{tikzcd}
\quad $\alpha(X) \circ u = \mathrm{id}_{X}$ \quad (unitarité)

\note{l'auteur porte (1.10) en tête et (1.9) au-dessous, les chiffres surchargés ;
c'est bien à (1.10) que renvoie la page 5}

On obtient un foncteur naturel
\[
  h : A \longrightarrow K(\varphi, \alpha, \lambda)
  \tag{1.11}
\]
par $Y \mapsto \bigl( f(Y),\ f(Y) \xrightarrow{\ f*\beta\ } \varphi f(Y) = fgf(Y) \bigr)$.

\marginal{vérifier compatibilités (1.9)}

\page{4}
\textbf{Théorème 1.12} (\uncertain{Beck}).
\struck{Supposons que} a) Si dans $A$ les $\mathrm{Ker}$ de doubles flèches
existent, et \struck{pour} \add{si} $f$ y commute $(\dagger)$, alors 1.11 est
essentiellement surjectif. b) Pour que $h$ soit une équivalence
[i.e. soit \struck{de plus fidèle} \add{pleinement fidèle}], il faut et il
suffit que $f$ soit \emph{conservatif}, et que pour toute double flèche $(u,v)$
de $A$ telle que $\mathrm{Ker}(f(u), f(v))$ existe, $\mathrm{Ker}(u,v)$ existe
et $f$ y commute.

\marginal{$(\dagger)$ il suffit de regarder les doubles flèches $u, v$ de $A$
telles que $\mathrm{Ker}(f(u), f(v))$ existe \quad (condition C)}

\marginal{\ill{} auxiliaires a), on établissant que $h$ est une
\uncertain{équivalence} \ill{}}
\note{seconde note marginale, écrite en biais et en grande partie illisible}

\subsection*{2) La réciproque : de $(\varphi,\alpha,\lambda)$ au couple adjoint}

2) Partons d'une catégorie $B$, avec un foncteur
\[
  \varphi : B \longrightarrow B
  \tag{2.1}
\]
\struck{satisfaisant les conditions (1.5),} et des liaisons
\[
  \varphi \xrightarrow{\ \alpha\ } \mathrm{id}_{B} , \qquad
  \varphi \xrightarrow{\ \lambda\ } \varphi^{2}
  \tag{2.2}
\]
satisfaisant les conditions (1.5), (1.6). On peut alors (= sans ces
restrictions (1.5) et (1.6)\,!) définir une catégorie
\[
  A = K(\varphi, \alpha, \lambda)
  \tag{2.3}
\]
comme ci-dessus. Soit
\[
  A \xrightarrow{\ f\ } B
  \tag{2.4}
\]
le foncteur \struck{d'oubli} canonique $(X, u) \mapsto X$.

\struck{Question. Quand $f$ admet-il un adjoint à droite $g$, de telle façon
que \ill{} que $fg$ \ill{}}

On définit un foncteur
\[
  g : B \longrightarrow A
  \tag{2.5}
\]
par
\[
  g(X) = \bigl( \varphi(X),\ \varphi(X) \xrightarrow{\ \lambda(X)\ } \varphi^{2}(X) \bigr)
  \tag{2.6}
\]
[la condition (1.6) assurant qu'on tombe bien dans $K(\varphi,\alpha,\lambda)$\,!].

\struck{$\mathrm{Hom}_{A}((X,u), g(Y)) \simeq \mathrm{Hom}_{B}(f(X,u), Y)$ \ill{}}

\[
  fg = \varphi
  \tag{2.7}
\]

\page{5}
Considérons donc
\[
  fg \xrightarrow{\ \alpha\ } \mathrm{id}
  \tag{2.8}
\]
ce qui fournit un homomorphisme d'adjonction
\[
  \alpha' : \mathrm{Hom}_{A}(X, g(Y)) \longrightarrow \mathrm{Hom}_{B}(f(X), Y) ,
  \qquad w \mapsto \alpha \circ f(w)
  \tag{2.9}
\]
\struck{$v \mapsto \alpha \circ f(v)$} \quad (2.9$'$)

On va définir aussi
\[
  \beta : \mathrm{id}_{K} \longrightarrow gf
  \tag{2.10}
\]
en prenant, pour $Y = (X, u)$ dans $K$,

\begin{tikzcd}
X \arrow[r, "u"] \arrow[d, "u"'] & \varphi(X) \arrow[d, dashed, "\varphi(u)"] \\
\varphi(X) \arrow[r, "\lambda(X)"'] & \varphi^{2}(X)
\end{tikzcd}

\note{la colonne de gauche est $Y$, celle de droite $gf(Y)$}

la commutativité étant assurée par la condition (1.10) ; on a donc bien un
morphisme de $K$. On a aussi un homomorphisme d'adjonction en sens inverse
\[
  \beta' : \mathrm{Hom}_{B}(f(X), Y) \longrightarrow \mathrm{Hom}_{A}(X, g(Y)) ,
  \qquad w \mapsto g(w) \circ \beta(X)
  \tag{2.11}
\]
On vérifie que ces homomorphismes sont inverses l'un de l'autre. En effet
$\beta' \alpha' = \mathrm{id}$ résulte formellement de (1.5), et
$\alpha' \beta' = \mathrm{id}$ résulte de (1.9).

\struck{\ill{}}
\note{suit un calcul de six lignes entièrement biffé}

Donc $f$ et $g$ sont adjoints l'un de l'autre, et retrouvant bien sûr
$(\varphi, \alpha, \lambda)$ via \uncertain{le procédé} du n\textsuperscript{o}~1.

\page{6}
\textbf{Conclusion.} La donnée (pour $B$ fixé) d'une catégorie $A$, et d'un
couple de foncteurs adjoints $f, g$, $A \rightleftarrows B$, avec $f$
conservatif et \struck{commutant aux noyaux de doubles flèches} satisfaisant à
la condition C de 1.12, est essentiellement équivalente à la donnée d'un
foncteur $g$ dans $B$, avec les données de Godement $\alpha$, $\lambda$.

\struck{Corollaire. $f$ pleinement fidèle $\Leftrightarrow$ $\alpha * \varphi$
\ill{} $\Leftrightarrow$ $\lambda$ isom.}

\emph{Compléments sur propriétés d'exactitude.}

a) Pour tout type de $\varprojlim$ qui sont représentables dans $B$, les
$\varprojlim$ de ce type sont représentables dans $A$, et le foncteur $f$ y
commute ; \struck{si et seulement si $f$ y commute.}

b) Pour tout type de $\varprojlim$ qui est représentable dans $B$, \ill{} est
représentable dans $A$, et $f$ y commute.

c) Pour toute propriété d'exactitude de $B$, \struck{faisant intervenir} des
types de commutation de $\varprojlim_{J}$ \struck{de type I} avec des types de
$\varinjlim$ ou quelconques dans $B$, \ill{} la même propriété d'exactitude est
vraie dans $A$.
\note{ligne très surchargée ; la lecture des indices reste incertaine}

\struck{d) Supposons que $B$ soit une catégorie \ill{}}

\textbf{Théorème.} Supposons $B$ un topos. Alors $A$ est un topos et $f$ est de
la forme $f_{0}^{*}$ pour un morphisme de topos $f_{0} : B \to A$ \add{accessible}
(i.e. conservatif) sss $\varphi$ \struck{commute aux \ill{}} est exact à gauche.
Pour que $f_{0}$ soit \emph{essentiel}, il faut et il suffit que $\varphi$
commute aux $\varprojlim$ quelconques, i.e. aux produits infinis. (On suppose
que $\varphi$ a un adjoint à gauche $\psi$.)

\subsection*{3) Le cas $B = \prod_{i} B_{i}$}

\page{7}
3) Supposons
\[
  B = \prod_{i \in I} B_{i}
  \tag{3.1}
\]
Donc $f$ est donné par une famille de
\[
  f_{i} : A \longrightarrow B_{i}
  \tag{3.2}
\]
et $g$ par une famille de
\[
  g_{i} : B_{i} \longrightarrow A ,
  \tag{3.3}
\]
\[
  g\bigl( (X_{i})_{i \in I} \bigr) = \prod_{i \in I} g_{i}(X_{i})
  \tag{3.4}
\]
Le composé $\varphi = fg$ est alors un foncteur
$\varphi : \prod_{i} B_{i} \to \prod_{i} B_{i}$, et est connu quand on connaît
les $\mathrm{pr}_{j}\, \varphi : \prod_{i} B_{i} \to B_{j}$. On trouve
\[
  \mathrm{pr}_{j}\, \varphi\bigl( (X_{i}) \bigr)
  = f_{j}\Bigl( \prod_{i \in I} g_{i}(X_{i}) \Bigr) .
  \tag{3.5}
\]
Soit
\[
  \varphi_{ji} \overset{\text{déf}}{=} f_{j} g_{i} : B_{i} \longrightarrow B_{j}
  \tag{3.6}
\]
et supposons que les $f_{j}$ commutent aux produits de type $I$ (p. ex. $f_{j}$
exacts ; ou $I$ fini). On a alors
\[
  \mathrm{pr}_{j}\, \varphi\bigl( (X_{i}) \bigr)
  = \prod_{i \in I} \varphi_{ji}(X_{i})
  \tag{3.6}
\]
\note{l'auteur écrit deux fois (3.6)}

et les données de $\alpha$, $\lambda$ deviennent celles de
\[
  \prod_{i} \varphi_{ki}(X_{i}) \xrightarrow{\ \bar{\alpha}_{k}\ } X_{k}
  \tag{3.7}
\]
\[
  \prod_{i} \varphi_{ki}(X_{i}) \xrightarrow{\ \bar{\lambda}_{k}\ }
  \prod_{i,j} \varphi_{kj} \varphi_{ji}(X_{i})
  \tag{3.8}
\]
i.e.
\[
  \prod_{\alpha} \varphi_{k\alpha}(X_{\alpha})
  \xrightarrow{\ \bar{\lambda}_{k,j,i}\ } \varphi_{kj} \varphi_{ji}(X_{i})
\]

\page{8}
qui s'expliciteront en termes des $\varphi_{ji} = f_{j} g_{i}$, et des
$\bar{\alpha}_{i}$, $\bar{\lambda}_{i}$ correspondants, ainsi
\[
  \left\{
  \begin{array}{l}
    \bar{\alpha}_{k} = \alpha_{k}(X_{k}) \circ \mathrm{pr}_{k} \\[2pt]
    \bar{\lambda}_{k,j,i} = \lambda_{kji}(X_{i}) \circ \mathrm{pr}_{i}
  \end{array}
  \right.
  \tag{3.9}
\]
\[
  \lambda_{kji} : \varphi_{ki} \longrightarrow \varphi_{kj} \circ \varphi_{ji}
  \tag{3.10}
\]

\begin{tikzcd}[column sep=small, row sep=small, nodes={font=\scriptsize}]
\varphi_{ki} \arrow[r, "\lambda_{kji}"] \arrow[d, no head] & \varphi_{kj}\varphi_{ji} \arrow[d, no head] \\
f_{k} g_{i} & f_{k}\, g_{j} f_{j}\, g_{i}
\end{tikzcd}

\note{la flèche provient du morphisme d'adjonction $\mathrm{id} \to g_{j} f_{j}$,
que l'auteur indique par une flèche verticale sous $g_{j} f_{j}$}

Les $(\varphi_{ji}, \lambda_{kji})$ sont reliées par des conditions exprimant
(1.5), (1.6), que nous n'expliciterons pas.

\marginal{pour $i = j = k$, $\lambda_{kji}$ n'est autre que $\lambda_{i}$}

[Le cas le plus intéressant est celui où les \struck{$\varphi$}
\add{$\alpha_{i}$} sont des isomorphismes (i.e. les $g_{i}$ pleinement
fidèles). Dans ce cas, \struck{\ill{}} on peut supposer
$\varphi_{i} = \mathrm{id}_{B_{i}}$,
$\alpha_{i} = \mathrm{id}_{\mathrm{id}_{B_{i}}}$, $\lambda_{i} = \mathrm{id}$.]

Les conditions (1.) \struck{les} dans les cas où deux des $i, j, k$ sont égaux,
avec $i = j$ ou $j = k$, \uncertain{ne montrent que} l'un des $\lambda_{kji}$
est l'identité, \ill{} de $\varphi_{kj} \varphi_{ji}$.
\note{ligne très surchargée ; le verbe reste douteux}

Il reste donc à examiner les cas (3.11) : a) \struck{$i$, $j$, $k$ deux à deux
différents} ; b) $i = k \neq j$, d'où
\[
  \mathrm{id}_{B_{i}} \longrightarrow \varphi_{ij} \circ \varphi_{ji}
  \qquad (i \neq j)
\]

\page{9}
Prenons le cas
\[
  B = B' \times B'' , \quad f = (f', f'') , \quad
  g(X', X'') = g'(X') \times g''(X'')
  \tag{3.12}
\]
Supposons $f' g' \simeq \mathrm{id}_{B'}$, $f'' g'' \simeq \mathrm{id}_{B''}$
(i.e. $g'$, $g''$ pleinement fidèles), donc $\varphi$ est \struck{connue}
\add{connu} quand on connaît
\[
  \varphi'' : B' \longrightarrow B'' , \qquad
  \varphi' : B'' \longrightarrow B'
  \tag{3.13}
\]
et la donnée des \struck{$\alpha$} $\lambda_{kji}$ devient celle de

\struck{$X' \times \varphi'(X'') \to X'$ , $\varphi''(X') \times X'' \to X''$}

\struck{$X' \to \varphi' \varphi''(X')$ , $X'' \to \varphi'' \varphi'(X'')$}

\note{deux blocs numérotés puis biffés en croix ; l'auteur y compte les cas
$(i,j,k)$ — « dont \struck{sept} 6 \struck{sont triviaux} » — avant de retenir
les seuls cas non triviaux ci-dessous}

il reste donc les données
\[
  \left\{
  \begin{array}{l}
    \mathrm{id}_{B'} \xrightarrow{\ \lambda'\ } \varphi' \varphi'' \\[2pt]
    \mathrm{id}_{B''} \xrightarrow{\ \lambda''\ } \varphi'' \varphi'
  \end{array}
  \right.
  \tag{3.14}
\]
Il faudrait \struck{\ill{}} expliciter dans ce cas les conditions (1.5) et (1.6),
qui portent sur
\[
  \varphi(X', X'') = \bigl( X' \times \varphi'(X'') ,\ \varphi''(X') \times X'' \bigr)
\]
et sur $\varphi^{2}(X', X'')$.
\note{le calcul est mené sur trois lignes puis biffé ; il aboutit aux relations
$\varphi'' * \lambda' = \lambda'' * \varphi''$ et
$\lambda' * \varphi' = \varphi' * \lambda''$, elles aussi biffées}

Explicitons la \struck{catégorie} $A$ en termes de $\varphi'$, $\varphi''$,
$\lambda'$, $\lambda''$. Elle a pour objets les quadruplets
$(X', X'', u', u'')$ où
\[
  X' \xrightarrow{\ u'\ } \varphi'(X'') , \qquad
  X'' \xrightarrow{\ u''\ } \varphi''(X')
  \tag{3.15}
\]

\page{10}
(définissant \struck{\ill{}}
$X' \xrightarrow{(\mathrm{id}_{X'},\, u')} X' \times \varphi'(X'')$,
$X'' \xrightarrow{(\mathrm{id}_{X''},\, u'')} X'' \times \varphi''(X')$, ce qui
exprime déjà (1.9)) soumis à la condition de commutativité de (1.10)
\[
  (X', X'') \xrightarrow{\ u\ }
  \bigl( X' \times \varphi'(X'') ,\ X'' \times \varphi''(X') \bigr)
  \rightrightarrows \varphi^{2}(X', X'')
\]
les deux flèches étant $\varphi(u)$ et $\lambda$, i.e. via $u'$, $u''$ d'une
part et via $\lambda'$, $\lambda''$ de l'autre. Il semble qu'on trouve la
commutativité des diagrammes :

(3.16)

\begin{tikzcd}
X' \arrow[r, "u'"] \arrow[rr, bend right=30, "\lambda'(X')"'] & \varphi'(X'') \arrow[r, "\varphi'(u'')"] & \varphi'\varphi''(X')
\end{tikzcd}

\begin{tikzcd}
X'' \arrow[r, "u''"] \arrow[rr, bend right=30, "\lambda''(X'')"'] & \varphi''(X') \arrow[r, "\varphi''(u')"] & \varphi''\varphi'(X'')
\end{tikzcd}

\textbf{Exemple.} Si $\varphi'' \varphi'$ et $\varphi' \varphi''$ sont des
foncteurs constants de valeurs un objet final de $B'$ resp. $B''$, les
conditions 3.16 ci-dessus sont vides, et on trouve pour $A$ la catégorie des
objets
\[
  \bigl( X', X'',\ X' \xrightarrow{\ u'\ } \varphi'(X'') ,\
  X'' \xrightarrow{\ u''\ } \varphi''(X') \bigr)
\]
$f'$ et $f''$ étant $\mapsto X'$ et $\mapsto X''$, et
\[
  \left\{
  \begin{array}{l}
    g'(X') = \bigl( X',\ \varphi''(X'),\
      X' \xrightarrow{\lambda'(X')} \varphi'\varphi''(X') ,\
      \varphi''(X') \xrightarrow{\mathrm{id}} \varphi''(X') \bigr) \\[3pt]
    g''(X'') = \bigl( \varphi'(X''),\ X'',\
      \varphi'(X'') \xrightarrow{\mathrm{id}} \varphi'(X'') ,\
      X'' \xrightarrow{\lambda''(X'')} \varphi''\varphi'(X'') \bigr)
  \end{array}
  \right.
  \tag{3.17}
\]
\[
  \left\{
  \begin{array}{l}
    \text{Image essentielle de } g' : u'' \text{ est un isomorphisme} \\[2pt]
    \text{Image essentielle de } g'' : u' \text{ est un isomorphisme}
  \end{array}
  \right.
  \tag{3.18}
\]

\page{11}
L'auteur écrit le morphisme d'adjonction $\beta'$ colonne par colonne :
\[
  \begin{array}{ccl}
    X & = & \bigl( X', X'',\ X' \xrightarrow{u'} \varphi'(X'') ,\
      X'' \xrightarrow{u''} \varphi''(X') \bigr) \\[4pt]
    \downarrow \beta' & & \\[4pt]
    g' f'(X) & = & \bigl( X',\ \varphi''(X'),\
      X' \xrightarrow{\lambda'(X')} \varphi'\varphi''(X') ,\
      \varphi''(X') \xrightarrow{\mathrm{id}} \varphi''(X') \bigr)
  \end{array}
  \tag{3.20}
\]
les composantes de $\beta'$ étant $\mathrm{id}$, $u''$, puis $\mathrm{id}$,
$\varphi'(u'')$ et $u''$, $\mathrm{id}$ ; et $\beta''$ analogue.
\marginal{adjonction}
\note{l'auteur passe de (3.18) à (3.20) sans écrire de (3.19)}

[Question. Quand $\beta'$ est-il un épimorphisme pour tout $X$ ? Cela
\struck{signifie} implique, moyennant des conditions d'exactitude
\uncertain{à droite} standard sur $\varphi'$, $\varphi''$, que
$u'' : X'' \to \varphi''(X')$ est un épimorphisme.]

L'intersection des images essentielles de $g'$ et de $g''$ est formée des
$(X', X'', u', u'')$ pour lesquels $u'$ et $u''$ sont deux isomorphismes.
\struck{Cette catégorie} Cette sous-catégorie, image commune de $g'$ et $g''$,
est équivalente à la sous-catégorie pleine de $B'$ formée des $X'$ tels que
$\lambda'(X')$ soit un isomorphisme.

(3.21) \quad Conditions équivalentes :

(i) $g'_{\mathrm{ess}}(B') \cap g''_{\mathrm{ess}}(B'')$ est une catégorie
\uncertain{ponctuelle} ;

(ii) la sous-catégorie pleine de $B'$ formée des $X'$ tels que
$\lambda'(X')$ soit un isomorphisme est \uncertain{ponctuelle} ;

(ii bis) idem pour $\lambda''$.

\section*{Résolutions standard et cohomologie relative (pages 13 à 20)}

\subsection*{Résolutions standard et théories relatives}

\page{13}
\emph{Résolutions standard et théories relatives.}

\begin{tikzcd}
C \arrow[r, bend left=20, "F"] & C' \arrow[l, bend left=20, "G"]
\end{tikzcd}
\qquad
$\mathrm{Hom}\bigl( F(X), X' \bigr) \simeq \mathrm{Hom}\bigl( X, G(X') \bigr)$

Posons $E = GF : C \to C$, d'où
\[
  \phi : \mathrm{id}_{C} = E^{0} \longrightarrow E , \qquad
  \lambda : E^{2} \longrightarrow E
\]
satisfaisant les relations de Godement, d'où un objet simplicial $\widehat{E}$
dans $\mathrm{Hom}(C, C)$ : $\widehat{E}_{n} = E^{n+1}$, et augmentation
$X \to \widehat{E}(X)$. (On devrait donner une interprétation naturelle par
théorie de la descente.)

On pose de même $E' = FG : C' \to C'$, d'où
\[
  \phi' : E' \longrightarrow \mathrm{id}_{C'} = E'^{0} , \qquad
  \lambda' : E' \longrightarrow E'^{2}
\]
satisfaisant aux relations de Godement, d'où un objet simplicial
$\widehat{E}'$ dans $\mathrm{Hom}(C', C')$ : $\widehat{E}'_{n} = E'^{n+1}$.

Notons les \struck{isomorphismes simpliciaux}
\begin{enumerate}
  \item $G(X') \simeq \widehat{E}\bigl( G(X') \bigr)$ \quad pour
        $X' \in \mathrm{Ob}\, C'$ ;
  \item $F(X) \simeq F\bigl( \widehat{E}(X) \bigr)$ \quad pour
        $X \in \mathrm{Ob}\, C$.
\end{enumerate}
La dernière \ill{} un foncteur $T : C \to D$ et un homomorphisme fonctoriel
$h : T \circ E \to T$ tel que $h \circ (T * \phi) = \mathrm{id}_{T}$ ; alors
$T \widehat{E}(X) \simeq 0$ pour tout $X \in \mathrm{Ob}\, C$ [cf. Godement].
Les formules (i) se vérifient sans doute \struck{de même} par des formules
explicites.

\textbf{Application.} Supposons $C$, $C'$ abéliennes, $F$, $G$ additifs,
\struck{\ill{}}
\begin{enumerate}
  \item Si $F$ est exact et fidèle, alors $\widehat{E}(X)$ \ill{} une
        résolution de $X$ pour tout $X$.
  \item Si $\Gamma$ est un foncteur
\end{enumerate}
\note{la page s'arrête ici, en pleine phrase}

\subsection*{Cohomologie relative. 1. Caractérisation des foncteurs $A \mapsto \mathrm{Hom}(A,B)$}

\page{14}
\emph{Cohomologie relative} (un cas spécial).
\note{titre porté au crayon de couleur en tête de page}

1. \emph{Caractérisation des foncteurs $A \to \mathrm{Hom}(A, B)$.}

Soit $C$ une catégorie. Pour tout $B \in C$, soit $h_{B}$ le foncteur de $C$
dans les ensembles tel que
\[
  h_{B}(A) = \mathrm{Hom}(A, B) .
\]
$h_{B}$ est un foncteur, contravariant en $B$, à valeurs dans la catégorie des
foncteurs $C \to E$, et on a
\[
  \alpha : \mathrm{Hom}(B, B') \xrightarrow{\ \sim\ }
  \mathrm{Hom}_{C}(h_{B}, h_{B'})
\]
[en effet, on obtient le foncteur inverse $\beta$ en associant à
$u \in \mathrm{Hom}_{C}(h_{B}, h_{B'})$ l'élément $\beta(u)$ de
$\mathrm{Hom}_{C}(B, B')$ image par $u_{B}$ de l'élément neutre de
$h_{B}(B) = \mathrm{Hom}_{C}(B, B)$ dans
$h_{B'}(B) = \mathrm{Hom}_{C}(B, B')$].

Si $C$ est additive, alors $h_{B}$ est additif [exact à gauche], et le foncteur
$B \mapsto h_{B}$ est additif (exact à gauche), et l'isomorphisme $\alpha$ est
valable en se limitant aux homomorphismes additifs. Ainsi $B \mapsto h_{B}$
permet d'« identifier » $C$ à une sous-catégorie de la catégorie
$\mathrm{Fonc.\,add.}(C^{\circ},\ $\uncertain{$\mathcal{Y}$}$)$.
\note{la catégorie but est notée d'une capitale cursive que la lecture ne
tranche pas ; elle revient plus bas et page 15}

Supposons que $C$ ait un générateur $U$, et \struck{\ill{}} satisfasse AB~3),
4), 5). \struck{Le foncteur} $h_{B}$ transforme limites inductives en limites
projectives, et dans ce cas quelconques. \ill{}

Si $h$, $h'$ sont deux foncteurs (contravariants) $C^{\circ} \to {}$%
\uncertain{$\mathcal{Y}$} \ill{} transformant limites inductives

\page{15}
en projectives, on \ill{} posant
\[
  M(h) = h(U) ,
\]
on a
\[
  \alpha' : \mathrm{Hom}(h, h') \xrightarrow{\ \sim\ }
  \mathrm{Hom}_{R^{\circ}}\bigl( M(h), M(h') \bigr)
\]
où $R = \mathrm{Hom}_{C}(U, U)$, $R^{\circ} = $ anneau inverse de $R$.

Pour le voir, on définit un homomorphisme
\[
  \beta : \mathrm{Hom}_{R^{\circ}}\bigl( M(h), M(h') \bigr)
  \longrightarrow \mathrm{Hom}(h, h')
\]
[détails laissés au lecteur]. Donc le foncteur $h \mapsto M(h)$ identifie la
catégorie des foncteurs \struck{\ill{}} envisagés à une sous-catégorie de
$\mathrm{Mod}_{R^{\circ}}$ ($R$-modules à droite).

On va définir un foncteur « inverse » $M \mapsto h_{M}$ de
$\mathrm{Mod}_{R^{\circ}}$ dans la catégorie de foncteurs envisagée ;
définissons le foncteur $T'' : \mathrm{Mod}_{R^{\circ}} \to C$, ou plus
simplement
\[
  T''(M) = M \otimes_{R} U .
\]
On va montrer que $T' T''$ est isomorphe à l'identité, donc $T''$ et $T'$ sont
des « équivalences ». On doit donc montrer que
\[
  \left\{
  \begin{array}{l}
    \mathrm{Hom}_{C}(U,\ M \otimes_{R} U) \simeq M \\[2pt]
    \text{si } R = \mathrm{Hom}_{C}(U, U) , \quad M \in \mathrm{Mod}_{R^{\circ}}
  \end{array}
  \right.
\]
ce qui est bien connu, à cause de \ill{}. Ainsi, les foncteurs
\struck{covariants} \add{contravariants} \add{exacts} à valeurs
$C \to {}$\uncertain{$\mathcal{Y}$}, transformant les limites inductives en
limites projectives, sont essentiellement les $h_{B}$.
\marginal{les $h_{B}$}

\subsection*{2. Le couple $(S, T)$ et les morphismes $\varphi$, $\psi$}

\page{16}
2. Soient $C$, $C'$ deux catégories \struck{abéliennes} \add{additives},
$S : C \to C'$ un foncteur additif de $C$ dans $C'$. On suppose que pour tout
\struck{$A \in C$} $X \in C'$, le foncteur
$A \mapsto \mathrm{Hom}_{C'}\bigl( S(A), X \bigr)$ est un foncteur du type
\struck{$\mathrm{Hom}(A,$} $h_{B}$ (i.e. $B$ dépend de $X$, mais de façon
fonctorielle), donc on aura
\[
  \mathrm{Hom}_{C'}\bigl( S(A), X \bigr) \simeq
  \mathrm{Hom}_{C}\bigl( A, T(X) \bigr)
  \tag{1}
\]
où $T : C' \to C$ est un foncteur covariant exact : [foncteur de $C'^{\circ}$
dans $C$]. Cette condition est \struck{\ill{}} satisfaite notamment si $C$ et
$C'$ vérifient AB~3), 4), 5), et $C$ a un générateur, et si de plus
\struck{\ill{}} $S$ est un foncteur exact à droite (donc
$A \mapsto \mathrm{Hom}_{C'}(S(A), X)$ est exact à gauche).

\marginal{il faudrait des conditions de représentabilité \ill{}}

En particulier, si $X = S(B)$, on aura
\[
  \mathrm{Hom}_{C}(A, B) \longrightarrow
  \mathrm{Hom}_{C'}\bigl( S(A), S(B) \bigr) \simeq
  \mathrm{Hom}_{C}\bigl( A, TS(B) \bigr)
  \tag{2}
\]
Donc, faisant $A = B$, on trouve un élément canonique
\[
  \varphi(A) \in \mathrm{Hom}_{C}\bigl( A, TS(A) \bigr) , \quad \text{i.e.}
  \tag{3}
\]
\[
  \varphi : \mathrm{id}_{C} \longrightarrow TS
  \tag{3 bis}
\]
D'autre part, on \struck{a} tire de (1), faisant $A = T(Y)$,
\[
  \mathrm{Hom}_{C'}\bigl( ST(Y), X \bigr) \simeq
  \mathrm{Hom}_{C}\bigl( T(Y), T(X) \bigr)
  \tag{4}
\]
et en particulier, faisant $X = Y$, on trouve un élément canonique
\[
  \psi(X) \in \mathrm{Hom}_{C'}\bigl( ST(X), X \bigr) , \quad \text{i.e.}
  \tag{5}
\]
\[
  \psi : ST \longrightarrow \mathrm{id}_{C'}
  \tag{5 bis}
\]

\page{17}
Considérons, pour $A \in C$, $\varphi(A) : A \to TS(A)$, d'où
$S(\varphi(A)) : S(A) \to STS(A)$ ; d'autre part, on a un morphisme
$(ST)(SA) \xrightarrow{\psi(SA)} SA$. Ceci dit, on a le composé

\begin{tikzcd}
S(A) \arrow[r, "S(\varphi(A))"] \arrow[rr, bend right=30, "\text{identité}"'] & S(TS(A)) = (ST)(S(A)) \arrow[r, "\psi(S(A))"] & S(A)
\end{tikzcd}

ou, en d'autres termes, $S(A)$ s'identifie par $S(\varphi(A))$ à un facteur
direct de $(ST)(S(A))$, la projection naturelle étant $\psi(S(A))$. De même,
on a le diagramme commutatif

\begin{tikzcd}
T(X) \arrow[r, "\varphi(T(X))"] \arrow[rr, bend right=30, "\text{identité}"'] & T(ST(X)) = (TS)(T(X)) \arrow[r, "T(\psi(X))"] & T(X)
\end{tikzcd}

qui identifie $T(X)$ à un facteur direct de $(TS)(T(X))$.

Pour que l'homomorphisme fonctoriel $\varphi : \mathrm{id}_{C} \to TS$ soit
injectif (i.e. les $\varphi(A) : A \to TS(A)$ injectifs), il faut et il suffit
que les homomorphismes
\[
  S(A,B) : \mathrm{Hom}_{C}(A,B) \longrightarrow
  \mathrm{Hom}_{C'}\bigl( S(A), S(B) \bigr)
\]
soient injectifs ; ce qui équivaut aussi à ce que le foncteur $S$ soit
\struck{\ill{}} « \uncertain{injectif} ». Cela résulte de la commutativité de

\begin{tikzcd}[column sep=small, nodes={font=\scriptsize}]
\mathrm{Hom}_{C}(A,B) \arrow[r] \arrow[rr, bend right=30] & \mathrm{Hom}_{C'}\bigl( S(A), S(B) \bigr) & \mathrm{Hom}_{C}\bigl( A, TS(B) \bigr) \arrow[l, "\simeq"']
\end{tikzcd}

\subsection*{3. Suites $S$-exactes, objets $S$-injectifs}

\page{18}
3) Une suite
\[
  0 \longrightarrow A' \longrightarrow A \longrightarrow A'' \longrightarrow 0
  \tag{1}
\]
dans $C$ est dite \emph{$S$-exacte} si la suite
\[
  0 \longrightarrow S(A') \longrightarrow S(A) \longrightarrow S(A'')
  \longrightarrow 0
  \tag{1 bis}
\]
est exacte et se scinde. Un objet $M \in C$ est dit \emph{$S$-injectif} si,
pour toute suite $S$-exacte (1), la suite
\[
  0 \longrightarrow \mathrm{Hom}(A'', M) \longrightarrow \mathrm{Hom}(A, M)
  \longrightarrow \mathrm{Hom}(A', M) \longrightarrow 0
  \tag{2}
\]
est exacte.

\textbf{Exemples.} (i) $M = T(X)$, car alors la suite (2) s'écrit
\[
  0 \longrightarrow \mathrm{Hom}\bigl( S(A''), X \bigr) \longrightarrow
  \mathrm{Hom}\bigl( S(A), X \bigr) \longrightarrow
  \mathrm{Hom}\bigl( S(A'), X \bigr) \longrightarrow 0
  \tag{2 bis}
\]
[D'ailleurs, si (1) est telle que pour tout objet $S$-injectif $M$
\struck{\ill{}} de la forme $M = T(X)$ la suite (2) est exacte, alors (1) est
$S$-exacte, car faisant \ill{} $X$ dans (1 bis) exacte ; (2 bis) sera exacte
pour \struck{\ill{}} $X = S(A')$.]
\note{le raisonnement entre crochets est deux fois repris et biffé}

(ii) Supposons $S$ exact \add{à droite} et « fidèle » \struck{en ce sens que}
$S(u) = 0 \Rightarrow u = 0$. Soit $A \in C$ ; on dit que la suite
\[
  \Sigma(A) : \quad 0 \longrightarrow A \longrightarrow TS(A)
  \longrightarrow Z(A) \longrightarrow 0
  \tag{3}
\]
(où $Z(X) = \mathrm{coker}\,\bigl( X \to TS(X) \bigr)$) est $S$-exacte. En
effet, elle se transforme par $S$

\page{19}
en
\[
  0 \longrightarrow S(A) \longrightarrow STS(A) \longrightarrow SZ(X)
  \longrightarrow 0
\]
elle est exacte car \ill{} $STS(A)$ et $SZ(X)$ puisque $S$ est exact à droite,
et \ill{} elle est scindée, ce qui \ill{} explicite une rétraction construite
plus haut de $S(A) \to S(TS)(A) = (ST)S(A)$.

[Soit $M \in C$ tel que, pour toute suite $S$-exacte de la forme $\Sigma(A)$,
la suite $\mathrm{Hom}\bigl( \Sigma(A), M \bigr)$ soit exacte :
\[
  0 \longrightarrow \mathrm{Hom}\bigl( Z(A), M \bigr) \longrightarrow
  \mathrm{Hom}\bigl( TS(A), M \bigr) \longrightarrow
  \mathrm{Hom}(A, M) \longrightarrow 0
\]
Alors $M$ est $S$-injectif. En effet, on fait $M = A$, et on trouve que
$A \to TS(A)$ identifie $A$ à un \emph{facteur direct} de $TS(A) = T(X)$
($X = S(A)$). Or $T(X)$ est $S$-injectif en vertu de (i), donc tout facteur
direct l'est, donc $A$ l'est. Ce qui \struck{$\Rightarrow$} donne que
$M$ est $S$-injectif $\Longleftrightarrow$ $M \to TS(M)$ \ill{} identifie $M$ à
un facteur direct (cf. les faits de nature \ill{} de la catégorie
Eilenberg--Moore).]

\textbf{Remarque.} \ill{} \struck{pour que le foncteur $S$ soit exact}
\add{lorsque $S$ exact} \struck{i.e. $\varphi$ est injectif} et que
$S(u) = 0$ implique $u = 0$, alors une suite $S$-exacte est exacte. En effet,
si $\Sigma : A \xrightarrow{u} Q \xrightarrow{v} R$ et $S(v) \circ S(u) = 0$,
donc $S(v \circ u) = 0$, alors
\note{la remarque est réécrite trois fois et fortement biffée ; seule la
version portée en interligne se lit d'un bout à l'autre}

\page{20}
\struck{\ill{}}
\note{le haut de la page — un diagramme
$0 \to \mathrm{Hom}(B,A') \to \mathrm{Hom}(B,A) \to \mathrm{Hom}(B,A'') \to 0$
au-dessus de sa suite image par $S$, et une demi-page de commentaire — est
biffé en croix}

\ldots $v u = 0$, et on aura
\[
  S\bigl( \mathrm{Ker}\, v / \mathrm{Im}\, u \bigr) =
  \mathrm{Ker}\, S(v) / \mathrm{Im}\, S(u)
\]
et le \struck{des} premier membre est nul si et seulement si
$\mathrm{Ker}\, v / \mathrm{Im}\, u$ est nul, i.e. une suite $\Sigma$ est
exacte si et seulement si $S\Sigma$ l'est.

Nous avons maintenant ce qu'il faut

[catégories de « suites exactes » $\Longleftrightarrow$ catégories
d'« \ill{} »]

pour faire une théorie « \uncertain{véritable} » des foncteurs dérivés
$C \to D$, où $D$ est une \emph{véritable} catégorie abélienne. De plus,
définition des foncteurs $C^{0} = TS = C$ :
\[
  C^{1}(A) = C^{0}(A) / \mathrm{Im}\, A , \qquad
  C^{i}(A) = C^{0}\Bigl( C^{i-1}(A) / \mathrm{Im}\bigl( C^{i-2}(A) \bigr) \Bigr)
  \quad \text{pour } i \geqslant 2
\]
\marginal{($S$-injectifs)}
qui donnent une $S$-résolution canonique (fonctorielle) de $A \in C$

\end{document}
